\section{结论与展望}
\label{sec:conclusion}

本文针对捕获前自由漂浮空间机械臂的关节空间轨迹跟踪问题，构建了基于光滑周期延迟反馈的预设时间控制框架。该框架采用"扰动观测--动力学补偿--误差线性化--周期延迟反馈"的四阶段结构：首先利用动量守恒约束得到等效关节空间动力学，再通过预设时间扰动观测器估计集总扰动，随后采用计算力矩补偿将误差动力学规范化为二阶可控积分链，最后在该线性误差通道中引入光滑周期延迟反馈。本文的理论贡献可概括为以下三点：

\noindent\textbf{(i) 理想条件下的预设时间精确跟踪（定理~\ref{thm:main_ideal}）。} 在连续时间、精确模型、无饱和和精确补偿条件下，采用观测--PDF 两阶段时序时，闭环跟踪误差可在预设时间上界 \(T_{\mathrm{pre}}=T_o+2h\) 后精确归零。该上界完全由设计参数决定，与初始误差大小无关。Phase I 有界性由命题~\ref{prop:phaseI_ISS}~的 ISS 分析保证。

\noindent\textbf{(ii) 扰动条件下的实际有界跟踪（定理~\ref{thm:practical_boundary}、命题~\ref{prop:approx_null}）。} 存在观测误差、模型近似、执行器限制和数值离散时，本文将各非理想因素统一建模为等效残余加速度扰动 \(\Delta_r(t)\)，并建立了从残余扰动上界 \(\bar\Delta_r\) 到跟踪误差的定量映射关系。在精确 PDF 零化条件下，周期边界误差满足 \(\|z(t_k)\|\le \Gamma_h\bar\Delta_r\)；在近似零化条件下，稳态误差上界为 \(\Gamma_h\bar\Delta_r/(1-\varepsilon_\Delta)\)。

\noindent\textbf{(iii) Gramian 型 PDF 增益设计与维度敏感性分析（定理~\ref{thm:pdf_gain}）。} 给出了基于可控 Gramian 的 PDF 增益显式构造，使误差延迟系统的单周期状态转移矩阵实现零化。在此基础上，揭示了 Gramian 条件数随状态维数的增长规律：对于 \(n=7\) 的 MIMO 系统（\(2n=14\) 维），条件数可达 \(10^5\)--\(10^6\) 量级，导致实际闭环退化为命题~\ref{prop:approx_null}~所述的近似零化情形。仿真结果证实了该维度退化效应，明确了 PDF Gramian 构造的有效适用范围。

后续工作包括以下五个方面：
\begin{enumerate}
  \item 实现基于空间 CRBA 的完整 \(6+n\) 维浮基递推动力学仿真，并与名义模型进行交叉验证，以评估 Schur 补约化在不同工作区域内的近似精度。
  \item 针对 PDF Gramian 的病态性问题（在 \(n=7\) 时条件数可达 \(10^5\!\)--\(10^6\)），研究低灵敏度增益设计方法，如基于 \(\mathcal H_2/\mathcal H_\infty\) 优化的 Gramian 整形或结构化的延迟反馈参数化。
  \item 将末端任务空间轨迹生成、广义雅可比奇异规避、关节限位和动力学约束纳入统一控制框架，使 PDF 控制器在任务空间层面具备完整的工程适用性。
  \item 针对采样、测量噪声和执行器饱和等非理想因素，建立离散时间实际固定时间有界性理论，给出采样步长与稳态误差之间的显式关系。
  \item 在硬件在环（HIL）平台或气浮台上开展实验验证，检验所提控制结构在真实自由漂浮条件下的实际跟踪性能。
\end{enumerate}